\documentclass{article}
\usepackage[utf8]{inputenc}
\usepackage{amsmath}
\usepackage{geometry}
\geometry{margin=2cm}
\begin{document}

\section*{Resolução Completa da Questão 134 - ENEM 2024 - Ciências da Natureza}

\subsection*{Interpretação e Estratégia}
A questão aborda a limitação imposta pelo exoesqueleto rígido dos crustáceos para o seu crescimento e questiona como essa limitação é superada.

\textbf{Termos-Chave}
\begin{itemize}
  \item \textbf{Exoesqueleto}: revestimento externo rígido que protege e suporta o organismo.
  \item \textbf{Muda (Ecdise)}: processo de troca do exoesqueleto antigo por um novo, permitindo o crescimento.
\end{itemize}

\textbf{O que é pedido?} Identificar o mecanismo biológico utilizado pelos crustáceos para superar a limitação imposta pelo exoesqueleto rígido, possibilitando seu crescimento.

\textbf{Estratégia}: Compreender que o exoesqueleto rígido não permite crescimento contínuo, então o crustáceo realiza mudas periódicas para substituir esse revestimento por um maior. Avaliar as alternativas e escolher a que descreve corretamente este processo.

\subsection*{Conteúdo e Mini Revisão}
\begin{center}
\begin{tabular}{|p{4cm}|p{5cm}|p{5cm}|}
\hline
\textbf{Conceito} & \textbf{Idea-chave} & \textbf{Aplicação} \\
\hline
Exoesqueleto & Quitina e sais calcários; revestimento rígido & Limita crescimento por ser fixo e duro \\
\hline
Muda (Ecdise) & Troca do exoesqueleto para crescimento & Permite crescimento ao formar um novo exoesqueleto maior \\
\hline
Estrutura e crescimento & Exoesqueleto não cresce com o animal & Substituição periódica necessária para o crescimento \\
\hline
\end{tabular}
\end{center}

\textbf{Teoria:}
\begin{itemize}
  \item Crustáceos possuem exoesqueleto rígido que protege o corpo, mas não se expande.
  \item Para crescer, o crustáceo deve trocar esse revestimento periodicamente (muda ou ecdise).
  \item O exoesqueleto antigo é abandonado e um novo, maior, cresce por baixo, permitindo o aumento do tamanho corporal.
\end{itemize}

\subsection*{Resolução e "Pulo do Gato"}
\begin{enumerate}
  \item O exoesqueleto rígido limita o crescimento dos crustáceos.
  \item O crustáceo realiza a muda, trocando o exoesqueleto antigo por um novo maior.
  \item A alternativa D descreve corretamente esse processo.
  \item Portanto, a resposta correta é a alternativa D.
\end{enumerate}

\textbf{Pulo do Gato:} Lembre que o exoesqueleto é rígido e não cresce; o crustáceo cresce trocando o exoesqueleto na muda!

\subsection*{Análise dos Distratores}
\begin{itemize}
  \item \textbf{A}: incorreta, pois exoesqueleto está presente em todas as fases.
  \item \textbf{B}: errada, pois não há fragmentação para crescimento.
  \item \textbf{C}: equivocado, exoesqueleto rígido não molda continuamente.
  \item \textbf{E}: errado, composição não varia para crescimento.
\end{itemize}

\subsection*{Código da Questão}
CN_BIOLOGIA_MECANISMOS_001

%% Fim do documento completo

\end{document}